\chapter{Заключение}
\label{sec:conclusion}

\section{Основные результаты}

В работе спроектирован и формализован воспроизводимый пайплайн оценки
мультимодальных LLM-систем на основе LLM-as-a-Judge.

\begin{enumerate}
    \item Проведён обзор существующих подходов к формализованной оценке
        LLM (глава~\ref{sec:related}): бенчмарки и датасеты,
        LLM-as-a-Judge, доменно-специфические фреймворки, требования к
        воспроизводимости и известные ограничения (смещения CALM,
        чувствительность к промпту, контаминация данных).
        На основе обзора выбраны две опорные идеи -- PRD~\cite{li2023prd}
        и multi-prompt evaluation~\cite{mizrahi2024stateofwhatart}.
    \item Формализованы метрики оценки (глава~\ref{sec:method}):
        поточечная $\mathrm{NPS}$/$\Delta\mathrm{NPS}$
        (\eqref{eq:nps}--\eqref{eq:delta-nps}) и попарная
        $\mathrm{SbS\text{-}NPS}$ с пороговым голосованием
        (\eqref{eq:vote}--\eqref{eq:sbs-nps}) и
        bootstrap-доверительным интервалом
        (раздел~\ref{sec:method:bootstrap}).
    \item Формализована гибридная схема LLM-as-a-Judge
        (раздел~\ref{sec:method:hybrid}): классическая (single-judge),
        multi-prompt (\eqref{eq:multiprompt-judges}) и consensus
        (\eqref{eq:consensus-judges}) как частные случаи одного
        механизма голосования с разным числом элементов $\set{J}_q$.
    \item Описан протокол воспроизводимости пайплайна
        (раздел~\ref{sec:method:reproducibility}): версионирование
        моделей и промптов, фиксация параметров запуска, логирование и
        повторные прогоны.
    \item Проведены эксперименты по схеме из четырёх этапов
        (раздел~\ref{sec:experiments:design}), сравнивающие
        классическую, multi-prompt и consensus схемы между собой и с
        асессорской разметкой по $\mathrm{SbS\text{-}NPS}$ и стоимости.
\end{enumerate}

\textbf{Итоговый ответ на гипотезу.}
На наборе компактных открытых моделей-судей (N=50, текстовый режим) ни multi-prompt, ни consensus-схема не показали статистически значимого роста согласия с асессорами по сравнению с single-judge (парный bootstrap, H2 не подтверждена, раздел~\ref{sec:results:agreement}). Consensus при $l=4$ значимо переоценивает $\mathrm{SbS\text{-}NPS}$ и отклоняется от асессоров (нарушение эквивалентности \eqref{eq:tost}). Единственный подтверждённый эффект -- рост робастности покрытия: переход от одиночного судьи к multi-prompt ($3\times$) повышает долю запросов с валидным вердиктом с $34\%$ до $100\%$ (раздел~\ref{sec:results:agreement}). Сила моделей-судей (verbosity bias), а не механизм агрегации, является узким местом.

\section{Научный и практический вклад}

\textbf{Научный вклад} -- единая формализация классической,
multi-prompt и consensus схем LLM-as-a-Judge как частных случаев одной
модели голосования (раздел~\ref{sec:method:hybrid}), что позволяет
сравнивать их между собой и с асессорской разметкой по единой метрике
$\mathrm{SbS\text{-}NPS}$ с доверительным интервалом, а не по
разнородным метрикам, специфичным для отдельных фреймворков
(глава~\ref{sec:related}).

\textbf{Практический вклад} -- воспроизводимый пайплайн
(раздел~\ref{sec:method:reproducibility}), позволяющий регулярно (на
каждое изменение системы или конкурента) автоматически получать
сравнение $S$ и $C$ с количественной оценкой надёжности результата, без
обязательного привлечения асессоров на каждой итерации.
По модели стоимости (раздел~\ref{sec:experiments:cost}) даже самая
дорогая из рассмотренных автоматических схем (consensus при
$l \cdot m = 9$--$12$) обходится примерно на порядок дешевле полной
асессорской разметки той же корзины, что и делает регулярные прогоны
экономически осмысленными.

\section{Ограничения}

Результаты работы следует трактовать с учётом ряда ограничений.

\begin{enumerate}
    \item \textbf{Объём выборки.} Основной текстовый эксперимент проведён
        при $N = 50$ на test-сплите (MDE $\approx 20$~п.п.), что делает
        TOST с $\delta = 10$~п.п. недостижимым для этого набора данных.
        Мультимодальный ablation study проведён при $N = 200$
        (раздел~\ref{sec:results:multimodal-ablation}), что обеспечивает
        MDE $\approx 10$~п.п. и делает TOST теоретически достижимым.
        Тем не менее TOST не достигнут ни при $N = 50$, ни при $N = 200$ --
        это свидетельствует о реальном расхождении между автоматической
        оценкой и асессорской разметкой, а не только о недостатке мощности.
    \item \textbf{Однородность судей.} Использованные модели-судьи
        принадлежат к близким семействам открытых моделей; результаты
        могут не переноситься на другие архитектуры, а самопредпочтение
        внутри одного семейства консенсус устраняет лишь частично
        (раздел~\ref{sec:related:hybrid-baselines}).
    \item \textbf{Мультимодальность сведена к скриншотам.} Визуальный
        вход судьи приводится к статичному скриншоту ответа
        (раздел~\ref{sec:experiments:pilot-multimodal}); живые
        интерактивные элементы, видео и динамические превью остаются за
        рамками и требуют отдельного исследования.
    \item \textbf{Остаточная стохастичность.} Даже при температуре $0$
        часть API недетерминирована; разброс
        $\mathrm{SbS\text{-}NPS}$ между повторными прогонами $R$
        (раздел~\ref{sec:method:reproducibility}) не равен нулю и должен
        учитываться при сравнении схем.
    \item \textbf{Порог $\theta = 0.7$ взят из литературы.} Значение
        обосновано (раздел~\ref{sec:method:sbs-nps}), но без
        индивидуальной калибровки под конкретную корзину может быть
        неоптимальным.
    \item \textbf{Межасессорская согласованность (IRR).} Коэффициент конкордации (Cohen~$\kappa$) между асессорами на golden set не измерен. Это ограничение не позволяет количественно обосновать порог $\delta = 10$~п.п. для TOST (раздел~\ref{sec:method:bootstrap:equiv}) через межасессорскую дисперсию; порог обоснован косвенно (раздел~\ref{sec:experiments:power}).
\end{enumerate}

\paragraph{Что подтверждено и что не подтверждено.}
Таблица~\ref{tab:hypotheses-summary} сводит результаты проверки гипотез.

\begin{table}[htbp]
\centering
\caption{Результаты проверки гипотез}
\label{tab:hypotheses-summary}
\begin{tabular}{p{5cm}cc}
\toprule
Гипотеза & Результат & Причина \\
\midrule
Multi-prompt $>$ classic по согласию & Не подтверждена & Недостаток мощности / verbosity bias \\
Consensus $>$ multi-prompt & Не подтверждена & Verbosity bias усиливается \\
Схемы воспроизводят направление асессоров & Подтверждена & Знак $\mathrm{SbS\text{-}NPS}$ совпадает \\
Агрегация повышает робастность покрытия & Подтверждена & $34\% \to 100\%$ валидных вердиктов \\
\bottomrule
\end{tabular}
\end{table}

\section{Направления дальнейшей работы}

\begin{itemize}
    \item Итеративная дискуссия между судьями (полная схема
        PRD~\cite{li2023prd}) вместо одноразового консенсуса
        \eqref{eq:consensus-judges} -- оценка прироста согласия
        относительно роста стоимости.
    \item Адаптивное перекрытие $\abs{\set{J}_q}$: увеличение числа
        судей только для запросов с неоднозначным голосованием (близким
        к порогу $\theta$ в \eqref{eq:verdict}), снижающее среднюю
        стоимость по сравнению с фиксированным $l \cdot m$.
    \item Автоматическая калибровка порога $\theta$ и набора промптов
        $p_1, \dots, p_M$ по golden set
        (раздел~\ref{sec:method:data}) как часть пайплайна, а не как
        ручной шаг.
    \item Перенос формализации
        (раздел~\ref{sec:method:hybrid}) на мультимодальных судей,
        включая специфичные для изображений и видео типы смещений,
        отмеченные в MLLM-as-a-Judge~\cite{chen2024mllmjudge}.
    \item Регулярное обновление корзины $\set{Q}$ по аналогии с
        LiveBench~\cite{white2024livebench} для снижения риска
        контаминации и устаревания golden set.
    \item Интеграция пайплайна с практиками MLOps: автоматическое
        версионирование моделей, промптов и корзин, регистрация прогонов
        и метрик, регрессионные алерты на просадку
        $\mathrm{SbS\text{-}NPS}$ -- доведение протокола
        воспроизводимости (раздел~\ref{sec:method:reproducibility}) до
        непрерывного процесса оценки.
\end{itemize}
